% Propietario conservado: /Users/ruben/Documents/ChatGPT/jueces y controles/REVISION_ARGUMENTAL_APERTURAS_CIERRES_20260915/delivery/ARTICLE_V_ES_ARGUMENTAL_20260915/manuscrito/sections/62_medida_entropia_maxima.tex
% Resultado de máxima entropía recuperado del propietario de incidencia.
% Se desarrolla aquí su prueba variacional para hacerla local al artículo.
% No se atribuye novedad al resultado ni a su formulación convencional.
\subsection{Distribución estacionaria y caracterización variacional}
\label{v:v:subsec:parry-explicita}

Sea \(X_{\rm av}\) el espacio de sucesiones bilaterales sobre
\(\{\mathrm{ar},\mathrm{vol}\}\) cuyos pares consecutivos tienen
incidencia positiva en \(F_{\rm av}\). Es el envolvente de memoria uno
del calendario de modos: conserva sus pares admisibles, mientras permite
todas las concatenaciones compatibles con esos pares. Su medida
estacionaria pondera ese espacio de historias; la frecuencia cronológica
de la órbita \((\Sigma,\Pi,\Pi)\) sigue siendo, por separado,
\((1/3,2/3)\).

\begin{proposition}[Caracterización explícita de la medida de máxima entropía]
\label{v:v:prop:parry-explicita}
Con \(r=(1,\varphi)\), la matriz estocástica y su distribución
estacionaria son
\begin{equation}
 P_{ij}=\frac{(F_{\rm av})_{ij}r_j}{\varphi r_i},\qquad
 P=\begin{pmatrix}0&1\\\varphi^{-2}&\varphi^{-1}\end{pmatrix},
 \qquad
 \pi_{\rm av}=\frac{(1,\varphi^2)}{1+\varphi^2}.
 \label{v:v:eq:parry-transicion-estacionaria}
\end{equation}
La medida estacionaria de Markov \(\mu\) determinada por estos datos
tiene, para cada cilindro admisible de \(n\) transiciones,
\begin{equation}
 \mu[i_0\ldots i_n]
 =\frac{r_{i_0}r_{i_n}}{(1+\varphi^2)\varphi^n}.
 \label{v:v:eq:parry-cilindros}
\end{equation}
Es la única medida estacionaria sobre \(X_{\rm av}\) con tasa de
entropía máxima, y esa tasa es \(\log\varphi\).
\end{proposition}

\begin{proof}
La identidad \(\varphi^2=\varphi+1\) implica
\(\varphi^{-2}+\varphi^{-1}=1\), por lo que las filas de \(P\)
suman uno. La sustitución directa verifica
\(\pi_{\rm av}P=\pi_{\rm av}\). Las dos transiciones cruzadas
tienen probabilidad positiva; la cadena es irreducible, y la retención
en \(\mathrm{vol}\) la hace aperiódica. La distribución estacionaria
se obtiene también resolviendo las dos ecuaciones lineales y la
normalización, que dan la solución única impresa. La multiplicación de
\(\pi_{i_0}\prod_{k=0}^{n-1}P_{i_k i_{k+1}}\) cancela los factores
intermedios \(r_{i_k}\) y demuestra \eqref{v:v:eq:parry-cilindros}.

Para la maximalidad, sea \(\nu\) cualquier medida estacionaria
sobre \(X_{\rm av}\) y sea
\(q(j\mid\mathcal P)\) la distribución condicional de \(X_1\)
dado el pasado completo \(\mathcal P=(\ldots,X_{-1},X_0)\).
Su soporte está contenido en el de la fila \(P_{X_0,\cdot}\).
Para distribuciones finitas, la divergencia relativa
\(D(q\Vert p)=\sum_jq_j\log(q_j/p_j)\) es no negativa, y se
anula exactamente cuando \(q=p\). Esta propiedad se obtiene de
la concavidad del logaritmo, con la convención \(0\log0=0\).
En los pares admisibles,
\[
 \log P_{ij}=\log r_j-\log r_i-\log\varphi.
\]
La estacionariedad cancela la esperanza de
\(\log r_{X_1}-\log r_{X_0}\). Por tanto,
\begin{equation}
 h(\nu)=H_\nu(X_1\mid\mathcal P)
 =\log\varphi-
 \mathbb E_\nu D\!\left(q(\cdot\mid\mathcal P)
                  \Vert P_{X_0,\cdot}\right)
 \leq\log\varphi.
 \label{v:v:eq:parry-deficit-entropia}
\end{equation}
La identidad entre tasa y entropía condicional del pasado se obtiene
mediante la regla de la cadena sobre bloques finitos y el límite
decreciente de sus entropías condicionales; el alfabeto finito garantiza
que todas ellas estén acotadas entre cero y \(\log2\).

La medida \(\mu\) construida es de Markov con transición \(P\),
de modo que el término de divergencia se anula y
\(h(\mu)=\log\varphi\). Si otra medida estacionaria alcanza
esa tasa, la divergencia se anula casi seguramente: su distribución
condicional del futuro depende sólo de \(X_0\) y coincide con \(P\).
Es, por tanto, la misma cadena estacionaria, cuya distribución inicial
única es \(\pi_{\rm av}\). Esto demuestra la unicidad.
\end{proof}

La fórmula \eqref{v:v:eq:parry-deficit-entropia} precisa qué se maximiza:
la entropía por transición del envolvente simbólico declarado. La razón
\(\mu_{\rm ar}/\mu_{\rm vol}=\varphi^{-2}\) procede de sus pesos
estacionarios. La órbita cronológica periódica posee tasa de entropía
nula y mantiene su propio reparto temporal; ambos objetos permanecen
definidos con sus respectivos dominios.
